\section{Introduction}
A calibrated SOFR model can assign a variance to every future policy meeting. Which of those numbers are determined by derivative prices, and which are selected by the model? This question has two parts. Prices may fail to distinguish different allocations of meeting risk. Conversely, a chosen factor specification may rule out allocations that the prices would permit. We characterize both restrictions, then examine how much of the distinction can be measured in public option settlements.

The first result is an exact price-equivalence statement. In a Gaussian rate model with scheduled jumps and deterministic continuous volatility, suppose a family of cash claims observes the forward curve only after a group of meetings. The time-zero value of every fixed cash payoff in that family depends on the earlier meeting and continuous variances only through their total. The statement includes American cash exercise, even when exercise decisions may use the full earlier history. Futures accruing after the group retain one additional quantity: their first time-weighted variance moment. Consequently, with three or more earlier meetings, even complete strike surfaces and these futures generally leave a continuum of indistinguishable allocations. With no continuous volatility and a zero initial curve, three equally spaced dates with variances $(200,300,200)$ and $(300,100,300)$ bp$^2$ give the same prices throughout the stated contract family. The middle meeting's standard deviation nevertheless changes from $17.3$ to $10$ bp.

The second result determines restrictions imposed by the factor model. For the scheduled-rate construction of \citet{brace2024sofr}, we characterize all event-variance vectors attainable under arbitrary deterministic scale recalibration. The set is a finitely generated cone: subdividing the scale intervals beyond the meeting dates adds no attainable vectors. When each factor gives all future meetings the same loading, its complete restrictions are consecutive-meeting variance inequalities controlled by the largest decay rate. We then derive the conditional event-variance map with square-root stochastic volatility, including the random Heath--Jarrow--Morton (HJM) drift and its covariance with the diffusion. We determine its support, the closed set of possible meeting-variance vectors, and hence all affine equalities and the covariance rank. Applied to the two published parameter specifications, the result gives support dimension at most two for the constant specification and at most three for the time-dependent specification. Factor count alone would not reveal this difference.

The third result concerns the observation calendar and price precision. In the Gaussian model with deterministic continuous volatility, an ideal cash option expiring before its futures accrual begins reveals accumulated meeting variance plus weighted continuous variance. We give an exact rank criterion for separating the two, allowing several background-variance parameters and known maturity-dependent loadings. Premium intervals then produce sharp conditional ranges for individual meetings and group totals. Two suitably chosen cash calls already identify the exact price parameters when the initial curve is known. The calendar criterion and optimization are elementary once the pricing map is established. Their role is to show when a reported meeting allocation is unique, how background flexibility changes that answer, and whether uniqueness survives realistic price precision.

These results lead to different tests. The first identifies an ambiguity that additional post-meeting strikes and exercise opportunities cannot remove. The second identifies restrictions that no scale calibration of a specified factor model can evade. The third determines which earlier expiries and background exposures supply the missing information. Sections~\ref{sec:prices}--\ref{sec:precision} develop the price and calendar results; Section~\ref{sec:restrictions} treats factor restrictions.

\subsection{Relation to existing work}
Extracting announcement risk from options is established practice. \citet{patell1979anticipated} documents option-implied variance changes around anticipated earnings releases. \citet{dubinsky2018option} separates scheduled earnings uncertainty from ordinary volatility using option prices; that article subsumes the earlier two-author working paper. \citet{leung2014accounting} studies scheduled earnings jumps and American exercise. For interest rates, \citet{beber2006effect} examines changes in option-implied Treasury-price distributions around macroeconomic announcements, while \citet{wright2020event} uses short-dated options to separate event-day and ordinary-day risk. None of these makes the principle of subtracting background variance a new contribution here.

The closest monetary-policy measurement study is \citet{bauer2021market}, which constructs a model-free short-rate uncertainty measure from Eurodollar futures and options and studies its changes around FOMC announcements. Its Appendix C explicitly identifies a difficulty with an ex ante meeting-variance estimator: successive quarterly contracts need not span the same meetings, while monthly expiries have less historical data and liquidity. Our criterion addresses the cross-sectional version of that problem. It states exactly which meeting cutoffs and background exposures identify a particular future allocation, allows several background cells and different underlying windows, and distinguishes rank from useful resolution at finite price precision. It is complementary to their event-study evidence about how uncertainty changes through time.

\citet{zhong2026spanning} studies identification when a flexible background option surface can absorb scheduled-event risk. Our stronger price-equivalence result concerns an entire post-cutoff family of interest-rate cash claims, including exercise decisions using earlier information, and explains why futures preserve one more variance moment than discounted cash claims. This differs from showing that one variance estimator or one implied-volatility surface decomposition is nonunique.

Policy-date term-structure models provide the economic and pricing setting. \citet{piazzesi2005bond} models the relation between Federal Reserve policy and bond yields; \citet{kim2014jumps} treats yield jumps at known dates. \citet{gellert2021short} distinguishes predictable short-rate steps from continuously evolving expectations in the forward curve. \citet{backwell2021expected} separates scheduled and unscheduled overnight-rate jumps and studies their effects on backward-looking contracts. The HJM drift restriction \citep{heath1992bond} and the general treatment of predictable discontinuities in \citet{fontana2018term,fontana2022term} supply the consistency framework. We use explicit specializations of that framework to study the inverse pricing problem.

For compounded overnight rates, \citet{lyashenko2019looking} develops a joint framework for forward- and backward-looking rates, \citet{heitfield2019inferring} studies term-rate inference from SOFR futures, and \citet{skov2021dynamic} derives dynamic SOFR futures models. \citet{itkin2024semi} treats SOFR-futures option valuation, including American exercise. These works explain why accrual, discounting, and exercise conventions belong in the pricing map. In our setting, cash discounting removes the first time-weighted variance moment from post-cutoff curve-payoff values, whereas futures expectations retain it. For futures-style exercise, the absence of an early-exercise premium is the established convex-submartingale result of \citet{lieu1990option,chen1993pricing}. We use it to determine exactly which meeting aggregates survive the change of contract convention.

The closest factor-model comparison is \citet{brace2024sofr}. That paper constructs scheduled short-rate dynamics with continuous fixed-maturity forward rates and calibrates stochastic volatility to futures and option smiles. Our contribution is to characterize restrictions implied by its loading and variance-state specification: those surviving arbitrary deterministic scale choices, and the exact stochastic support for fixed parameters. The square-root transition law itself is classical \citep{malham2008square}; the meeting-variance support and its active-column rank follow only after accounting for the full conditional variance kernel. This is also distinct from unspanned stochastic volatility in \citet{collindufresne2002bonds}, where volatility risk is not spanned by bonds within a model. Here we distinguish parameter allocations that give identical selected prices and determine the set of meeting variances permitted by a fixed specification.

Finally, sharp financial bounds from incomplete observations have a substantial optimization literature. \citet{bertsimas2002relation} obtains bounds on prices and moments by convex optimization; \citet{beiglbock2011model} develops model-independent bounds using martingale transport. Our bounds are conditional on the Gaussian observation map and the discount factors and means used in inversion. The linear program is standard. The specific content is the feasible set generated by meeting dates, background exposures, and premium intervals, together with the exact indistinguishable-allocation set for post-cutoff contracts.

\subsection{Data application and scope}
The application uses 4,100 selected strike observations in 555 CME SOFR option chains on 78 dates. The retained calendars never identify the complete meeting/background vector. A normal-mixture smile fit reduces the median maximum strike error to 0.161 bp, but the resulting ATM meeting-risk bounds remain sensitive to background flexibility. Same-expiry options on later underlying quarters add information about maturity loadings. A hump that returns to its initial level fails on three-year midcurves; allowing a free long-end level reduces their prediction error from 3.818 to 1.131 bp in standard deviation. Positive constant-background meeting bounds survive wider exercise and mean allowances, while flexible-background bounds still include zero. A discrete-meeting example shows why these ATM-proxy allocations need not equal actual outcome variances. A companion second-moment calculation instead exposes sensitivity to tail extrapolation. These findings favor compatible ranges and explicit model checks over a single calibrated meeting-volatility curve.

There are three separate mathematical settings. Sections~\ref{sec:prices}--\ref{sec:calendar} include an arbitrary deterministic initial curve and deterministic continuous volatility. Section~\ref{sec:aggregation} first proves aggregation in the pure-jump specialization and then extends it to the initial curve and continuous volatility; its explicit exercise formulas use the specialization. Section~\ref{sec:restrictions} studies the different scheduled-rate factor construction described above, conditional on the stated square-root solutions. The public-data calculation uses normal and mixture-based ATM proxies for listed American options, with separate numerical exercise and mean sensitivities. Their calibration ranges are conditional on the stated proxy observation maps. Section~\ref{sec:empirical} separates its pricing errors from its calendar restrictions, and Section~\ref{sec:uses} gives the resulting hedge and convexity applications.
